% secciones/04_infraestructura.tex
\section{Infraestructura e Integración en el Entorno Hospitalario}
La eficacia operativa de un robot autónomo no reside únicamente en sus capacidades intrínsecas, sino en su grado de integración con la infraestructura física y digital del hospital. Un robot que opera de forma aislada ve limitada su utilidad; por ello, es fundamental proyectar un ecosistema de soporte que facilite su autonomía y garantice su disponibilidad continua.

Uno de los pilares de esta infraestructura es el despliegue estratégico de estaciones de recarga y mantenimiento, comúnmente denominadas \textit{docking stations}. Estas áreas no solo sirven como puntos de abastecimiento energético automatizado, donde el robot se acopla de forma precisa para recargar sus baterías, sino que también actúan como nodos de diagnóstico donde el sistema puede realizar chequeos de integridad de sus sensores y actuadores. Es imperativo que estas zonas estén ubicadas en puntos que no obstruyan el flujo asistencial, contando con sistemas de ventilación adecuados y protocolos de seguridad contra incendios específicos para baterías de alta densidad energética.

La comunicación bidireccional con el edificio inteligente representa otra capa crítica de la infraestructura. El robot debe ser capaz de interactuar con el sistema de gestión del edificio (BMS) para solicitar la apertura de puertas automáticas o la llamada de ascensores. Esta integración permite que el robot se desplace entre diferentes plantas sin intervención humana, optimizando las rutas logísticas y reduciendo el tiempo de entrega de materiales críticos. Este flujo de datos debe estar securizado bajo estándares de ciberseguridad industrial, evitando cualquier vulnerabilidad que pueda comprometer la red hospitalaria.

A diferencia de otros sistemas de desinfección ambiental, este dispositivo se especializa exclusivamente en la esterilización interna de su compartimento de carga mediante luz UV-C, garantizando la asepsia de las muestras biológicas transportadas. Por motivos de seguridad operativa y cumplimiento normativo, la desinfección del chasis exterior y de las superficies de contacto se delega en un protocolo de mantenimiento manual ejecutado por el personal de limpieza del centro, utilizando agentes químicos validados. Este enfoque garantiza que no existan riesgos de exposición accidental a radiación ultravioleta para los pacientes en los pasillos, manteniendo al robot como un sistema de transporte seguro y especializado.

Para asegurar la fiabilidad de la arquitectura Edge-Cloud, la infraestructura digital del hospital debe evolucionar hacia una red de baja latencia. Se prescribe el despliegue de una red inalámbrica con soporte para Fast Roaming, permitiendo que el robot transite entre puntos de acceso sin pérdida de paquetes críticos. Este ancho de banda garantizado es vital no solo para la telemetría, sino para permitir que el servidor central reasigne misiones en tiempo real en función de la demanda asistencial
